\pagestyle{empty}
\tight
\begin{tblr}{width=\textwidth,colspec={|Q[c,m,1.9cm]|X[l,m]|},hlines,vlines,hspan=minimal,row{1}={bg=headblue},rowsep=1pt,colsep=3pt}
\SetCell{c,m}\hfil\logo\hfil & \textbf{POLITEKNIK NEGERI MALANG}\par\textbf{JURUSAN TEKNOLOGI INFORMASI}\par\textbf{PROGRAM STUDI: D4 TEKNIK INFORMATIKA} \\
\SetCell[c=2]{c} \textbf{RENCANA PEMBELAJARAN SEMESTER (RPS)} & \\
\end{tblr}
\begin{longtblr}{width=\textwidth,colspec={|Q[l,m,1.9cm]|X[0.85,l,m]|X[1.45,l,m]|X[0.75,l,m]|X[1.15,l,m]|},hlines,vlines,hspan=minimal,rowsep=1pt,colsep=2pt,row{1,3}={bg=headgray,font=\bfseries}}
MATA KULIAH & KODE & BOBOT (SKS)/jam & SEMESTER & TGL. PENYUSUNAN \\
Sistem Operasi & RTI252004 & 2 SKS/4 jam & 2 & 06 Februari 2026 \\
OTORISASI & Dosen Pengembang RPS & Koordinator RMK & \SetCell[c=2]{l} Ka PRODI &  \\
 & Dian Hanifudin Subhi, S.Kom., M.Kom. &  & \SetCell[c=2]{l}{\vspace{8mm}\par Dr. Ely Setyo Astuti, S.T., M.T.} &  \\
\SetCell[r=14]{l} Capaian Pembelajaran (CP) & \SetCell[c=4]{l,bg=headgray,font=\bfseries} Capaian Pembelajaran Lulusan Program Studi (CPL-Prodi) &  &  &  \\
 & CPL02 & \SetCell[c=3]{l} Mampu menerapkan prinsip rekayasa sistem dan perangkat lunak untuk menghasilkan solusi aplikasi multi-platform sesuai kebutuhan. &  &  \\
 & CPL05 & \SetCell[c=3]{l} Mampu menghasilkan solusi teknologi komputasi dan infrastruktur digital yang sesuai dengan kebutuhan organisasi dan pengguna. &  &  \\
 & CPL09 & \SetCell[c=3]{l} Mampu menjelaskan cara kerja sistem komputer dan menerapkan pendekatan komputasi untuk menyelesaikan masalah. &  &  \\
 & \SetCell[c=4]{l,bg=headgray,font=\bfseries} Capaian Pembelajaran mata kuliah (CPL-MK) &  &  &  \\
 & CPMK0206 & \SetCell[c=3]{l} Mahasiswa mampu merancang strategi manajemen sumber daya sistem aplikasi multi-platform berdasarkan karakteristik sistem operasi. &  &  \\
 & CPMK0208 & \SetCell[c=3]{l} Mahasiswa mampu memilih arsitektur computing system sesuai kebutuhan aplikasi multi-platform. &  &  \\
 & CPMK0502 & \SetCell[c=3]{l} Mahasiswa mampu menganalisis isu keamanan sistem operasi, jaringan, dan cloud pada konteks layanan digital. &  &  \\
 & CPMK0901 & \SetCell[c=3]{l} Mahasiswa mampu menjelaskan cara kerja sistem komputer dan komponen utamanya dalam penyelesaian masalah komputasi. &  &  \\
 & \SetCell[c=4]{l,bg=headgray,font=\bfseries} Kemampuan akhir tiap tahapan belajar (Sub-CPMK) &  &  &  \\
 & SCPMK0901-01301 & \SetCell[c=3]{l} Mahasiswa mampu menjelaskan konsep dasar sistem operasi, kernel, distribusi Linux, instalasi, partisi, dan sistem file. &  &  \\
 & SCPMK0206-01302 & \SetCell[c=3]{l} Mahasiswa mampu menerapkan perintah terminal, I/O, proses, Bash, memori, dan system call untuk administrasi Linux. &  &  \\
 & SCPMK0502-01303 & \SetCell[c=3]{l} Mahasiswa mampu menganalisis permission, ACL, user/group, sudo, dan risiko keamanan administrasi sistem operasi. &  &  \\
 & SCPMK0208-01304 & \SetCell[c=3]{l} Mahasiswa mampu mengonfigurasi layanan, aplikasi, backup, recovery, dan remastering Linux sesuai kebutuhan sistem. &  &  \\
\SetCell[r=2]{l} Penilaian & \SetCell[c=4]{l} *Nilai CPMK didapat dari agregasi nilai SUB-CPMK yang dibebankan &  &  &  \\
 & \SetCell[c=4]{l}{\tiny\begin{tblr}{width=\linewidth,colspec={|Q[c,m,0.1\linewidth]|Q[c,m,0.16\linewidth]|X[l,m]|X[l,m]|X[l,m]|X[l,m]|X[l,m]|X[l,m]|Q[c,m,0.08\linewidth]|},hlines,vlines,hspan=minimal,rowsep=1pt,colsep=0.7pt,row{1,2}={bg=headgray,font=\bfseries}}
Kode CPMK & Kode SCPMK & \SetCell[c=6]{c} Bobot per Bentuk Penilaian & & & & & & Total Bobot \\
 & & Tugas & Kuis 1 & UTS & Kuis 2 & PBL & UAS & \\
CPMK0901 & SCPMK0901-01301 & 4.50\%: konsep OS dan instalasi & 5\%: materi minggu 1--3 & 5\%: administrasi Linux dasar & & & & 14.50\% \\
CPMK0206 & SCPMK0206-01302 & 7.50\%: terminal, I/O, proses, Bash, memori & & 20\%: administrasi Linux dasar & & & & 27.50\% \\
CPMK0502 & SCPMK0502-01303 & 10.50\%: file, user/group, akses & & & 5\%: manajemen akses & & 10\%: validasi administrasi sistem & 25.50\% \\
CPMK0208 & SCPMK0208-01304 & 20\%: layanan, aplikasi, backup, recovery & & & & 5\%: remastering Linux & 7.50\%: validasi administrasi sistem & 32.50\% \\
\SetCell[c=8]{r}\textbf{Total Per Penilaian} & & & & & & & & \textbf{100\%} \\
\end{tblr}} &  &  &  \\
Diskripsi Singkat Mata Kuliah & \SetCell[c=4]{l} Mata kuliah Sistem Operasi memberikan pengalaman praktik untuk memahami cara kerja sistem operasi dan melakukan administrasi Linux dari instalasi sampai layanan, keamanan akses, backup, recovery, dan remastering. &  &  &  \\
Bahan Kajian / Materi Pembelajaran & \SetCell[c=4]{l}\begin{enumerate}\item Konsep sistem operasi dan instalasi Linux\item Terminal, I/O, proses, Bash, memori, dan system call\item Permission, ACL, user/group, layanan, aplikasi, backup, recovery\item Remastering sistem operasi berbasis Linux\end{enumerate} &  &  &  \\
\SetCell[r=2]{l} Pustaka & \SetCell[c=4]{l} \textbf{Utama:}\begin{enumerate}\item Silberschatz, A., Galvin, P. B., and Gagne, G. 2018. Operating System Concepts, 10th Edition. Wiley.\item Stallings, W. 2017. Operating Systems: Internals and Design Principles, 9th Edition. Pearson.\end{enumerate} &  &  &  \\
 & \SetCell[c=4]{l} \textbf{Pendukung:} Materi LMS, dokumentasi resmi perangkat lunak, dan jobsheet praktikum. &  &  &  \\
Media Pembelajaran & \SetCell[c=2]{l} \textbf{Software:} Linux, terminal, shell, text editor, package manager, virtualization tools. &  & \SetCell[c=2]{l} \textbf{Hardware:} Komputer/laptop, proyektor, media instalasi/VM. &  \\
Nama Dosen Pengampu & \SetCell[c=4]{l}\begin{enumerate}\item Dr. Yuri Ariyanto, S.Kom., M.Kom.\item Agung Nugroho Pramuditha, S.T., M.T.\item Dian Hanifudin Subhi, S.Kom., M.Kom.\item Muhammad Afif Hendrawan, S.Kom., M.T.\item Imam Fahrur Rozi, S.T., M.T.\end{enumerate} &  &  &  \\
Matakuliah Syarat & \SetCell[c=4]{l} - &  &  &  \\
\end{longtblr}
\begin{longtblr}{width=\textwidth,colspec={|Q[c,m,0.06\textwidth]|X[1.35,l,m]|X[1.08,l,m]|X[1.16,l,m]|X[0.7,l,m]|X[1.24,l,m]|X[1.06,l,m]|X[0.98,l,m]|Q[c,m,0.05\textwidth]|},hlines,vlines,hspan=minimal,rowhead=2,row{1,2}={bg=headgreen,font=\bfseries},rowsep=1pt,colsep=1.5pt}
Mgg.\par Ke & Kemampuan Akhir Yang Direncanakan (Sub-CP-MK) & Bahan kajian (Materi Pembelajaran) & Bentuk dan Metode Pembelajaran & Estimasi Waktu & Pengalaman Belajar Mahasiswa & Kriteria \& Bentuk Penilaian & Indikator Penilaian & Bobot Penilaian (\%) \\
(1) & (2) & (3) & (4) & (5) & (6) & (7) & (8) & (9) \\
1 & SCPMK0901-01301 & Pengenalan OS dan instalasi Linux. & Modalitas: Blended Learning. Bentuk: Praktikum. Metode: demonstrasi, praktik laboratorium Linux, problem based learning, dan diskusi. & 1 x 2 x 50' tatap muka; 1 x 2 x 50' tugas/praktik mandiri. & Mahasiswa mengerjakan aktivitas terarah untuk pengenalan os dan instalasi linux dan menunjukkan evidence hasil belajar. & Bentuk penilaian: Pengenalan OS dan instalasi Linux. Mengacu rubrik RTM. & Ketepatan konsep; kualitas artefak/praktik; validasi dan dokumentasi. & 2 \\
2 & SCPMK0206-01302 & Perangkat keras dan perintah dasar. & Modalitas: Blended Learning. Bentuk: Praktikum. Metode: demonstrasi, praktik laboratorium Linux, problem based learning, dan diskusi. & 1 x 2 x 50' tatap muka; 1 x 2 x 50' tugas/praktik mandiri. & Mahasiswa mengerjakan aktivitas terarah untuk perangkat keras dan perintah dasar dan menunjukkan evidence hasil belajar. & Bentuk penilaian: Perangkat keras dan perintah dasar. Mengacu rubrik RTM. & Ketepatan konsep; kualitas artefak/praktik; validasi dan dokumentasi. & 3 \\
3 & SCPMK0206-01302 & Dasar input/output. & Modalitas: Blended Learning. Bentuk: Praktikum. Metode: demonstrasi, praktik laboratorium Linux, problem based learning, dan diskusi. & 1 x 2 x 50' tatap muka; 1 x 2 x 50' tugas/praktik mandiri. & Mahasiswa mengerjakan aktivitas terarah untuk dasar input/output dan menunjukkan evidence hasil belajar. & Bentuk penilaian: Dasar input/output. Mengacu rubrik RTM. & Ketepatan konsep; kualitas artefak/praktik; validasi dan dokumentasi. & 2 \\
4 & SCPMK0901-01301, SCPMK0206-01302 & Kuis 1. & Modalitas: Blended Learning. Bentuk: Praktikum. Metode: demonstrasi, praktik laboratorium Linux, problem based learning, dan diskusi. & 1 x 2 x 50' tatap muka; 1 x 2 x 50' tugas/praktik mandiri. & Mahasiswa mengerjakan aktivitas terarah untuk kuis 1 dan menunjukkan evidence hasil belajar. & Bentuk penilaian: Kuis 1. Mengacu rubrik RTM. & Ketepatan konsep; kualitas artefak/praktik; validasi dan dokumentasi. & 5 \\
5 & SCPMK0502-01303 & Struktur direktori dan operasi file. & Modalitas: Blended Learning. Bentuk: Praktikum. Metode: demonstrasi, praktik laboratorium Linux, problem based learning, dan diskusi. & 1 x 2 x 50' tatap muka; 1 x 2 x 50' tugas/praktik mandiri. & Mahasiswa mengerjakan aktivitas terarah untuk struktur direktori dan operasi file dan menunjukkan evidence hasil belajar. & Bentuk penilaian: Struktur direktori dan operasi file. Mengacu rubrik RTM. & Ketepatan konsep; kualitas artefak/praktik; validasi dan dokumentasi. & 3 \\
6 & SCPMK0206-01302 & Manajemen proses. & Modalitas: Blended Learning. Bentuk: Praktikum. Metode: demonstrasi, praktik laboratorium Linux, problem based learning, dan diskusi. & 1 x 2 x 50' tatap muka; 1 x 2 x 50' tugas/praktik mandiri. & Mahasiswa mengerjakan aktivitas terarah untuk manajemen proses dan menunjukkan evidence hasil belajar. & Bentuk penilaian: Manajemen proses. Mengacu rubrik RTM. & Ketepatan konsep; kualitas artefak/praktik; validasi dan dokumentasi. & 2 \\
7 & SCPMK0206-01302 & Bash shell. & Modalitas: Blended Learning. Bentuk: Praktikum. Metode: demonstrasi, praktik laboratorium Linux, problem based learning, dan diskusi. & 1 x 2 x 50' tatap muka; 1 x 2 x 50' tugas/praktik mandiri. & Mahasiswa mengerjakan aktivitas terarah untuk bash shell dan menunjukkan evidence hasil belajar. & Bentuk penilaian: Bash shell. Mengacu rubrik RTM. & Ketepatan konsep; kualitas artefak/praktik; validasi dan dokumentasi. & 3 \\
8 & SCPMK0901-01301, SCPMK0206-01302 & UTS administrasi Linux dasar. & Modalitas: Blended Learning. Bentuk: Praktikum. Metode: demonstrasi, praktik laboratorium Linux, problem based learning, dan diskusi. & 1 x 2 x 50' tatap muka; 1 x 2 x 50' tugas/praktik mandiri. & Mahasiswa mengerjakan aktivitas terarah untuk uts administrasi linux dasar dan menunjukkan evidence hasil belajar. & Bentuk penilaian: UTS administrasi Linux dasar. Mengacu rubrik RTM. & Ketepatan konsep; kualitas artefak/praktik; validasi dan dokumentasi. & 20 \\
9 & SCPMK0206-01302 & Bash programming. & Modalitas: Blended Learning. Bentuk: Praktikum. Metode: demonstrasi, praktik laboratorium Linux, problem based learning, dan diskusi. & 1 x 2 x 50' tatap muka; 1 x 2 x 50' tugas/praktik mandiri. & Mahasiswa mengerjakan aktivitas terarah untuk bash programming dan menunjukkan evidence hasil belajar. & Bentuk penilaian: Bash programming. Mengacu rubrik RTM. & Ketepatan konsep; kualitas artefak/praktik; validasi dan dokumentasi. & 2 \\
10 & SCPMK0206-01302 & Manajemen memori dan system call. & Modalitas: Blended Learning. Bentuk: Praktikum. Metode: demonstrasi, praktik laboratorium Linux, problem based learning, dan diskusi. & 1 x 2 x 50' tatap muka; 1 x 2 x 50' tugas/praktik mandiri. & Mahasiswa mengerjakan aktivitas terarah untuk manajemen memori dan system call dan menunjukkan evidence hasil belajar. & Bentuk penilaian: Manajemen memori dan system call. Mengacu rubrik RTM. & Ketepatan konsep; kualitas artefak/praktik; validasi dan dokumentasi. & 3 \\
11 & SCPMK0502-01303 & Manajemen file dan user/group. & Modalitas: Blended Learning. Bentuk: Praktikum. Metode: demonstrasi, praktik laboratorium Linux, problem based learning, dan diskusi. & 1 x 2 x 50' tatap muka; 1 x 2 x 50' tugas/praktik mandiri. & Mahasiswa mengerjakan aktivitas terarah untuk manajemen file dan user/group dan menunjukkan evidence hasil belajar. & Bentuk penilaian: Manajemen file dan user/group. Mengacu rubrik RTM. & Ketepatan konsep; kualitas artefak/praktik; validasi dan dokumentasi. & 5 \\
12 & SCPMK0502-01303 & Kuis 2. & Modalitas: Blended Learning. Bentuk: Praktikum. Metode: demonstrasi, praktik laboratorium Linux, problem based learning, dan diskusi. & 1 x 2 x 50' tatap muka; 1 x 2 x 50' tugas/praktik mandiri. & Mahasiswa mengerjakan aktivitas terarah untuk kuis 2 dan menunjukkan evidence hasil belajar. & Bentuk penilaian: Kuis 2. Mengacu rubrik RTM. & Ketepatan konsep; kualitas artefak/praktik; validasi dan dokumentasi. & 5 \\
13 & SCPMK0208-01304 & Manajemen layanan. & Modalitas: Blended Learning. Bentuk: Praktikum. Metode: demonstrasi, praktik laboratorium Linux, problem based learning, dan diskusi. & 1 x 2 x 50' tatap muka; 1 x 2 x 50' tugas/praktik mandiri. & Mahasiswa mengerjakan aktivitas terarah untuk manajemen layanan dan menunjukkan evidence hasil belajar. & Bentuk penilaian: Manajemen layanan. Mengacu rubrik RTM. & Ketepatan konsep; kualitas artefak/praktik; validasi dan dokumentasi. & 5 \\
14 & SCPMK0208-01304 & Manajemen aplikasi. & Modalitas: Blended Learning. Bentuk: Praktikum. Metode: demonstrasi, praktik laboratorium Linux, problem based learning, dan diskusi. & 1 x 2 x 50' tatap muka; 1 x 2 x 50' tugas/praktik mandiri. & Mahasiswa mengerjakan aktivitas terarah untuk manajemen aplikasi dan menunjukkan evidence hasil belajar. & Bentuk penilaian: Manajemen aplikasi. Mengacu rubrik RTM. & Ketepatan konsep; kualitas artefak/praktik; validasi dan dokumentasi. & 5 \\
15 & SCPMK0208-01304 & Backup dan pemulihan sistem. & Modalitas: Blended Learning. Bentuk: Praktikum. Metode: demonstrasi, praktik laboratorium Linux, problem based learning, dan diskusi. & 1 x 2 x 50' tatap muka; 1 x 2 x 50' tugas/praktik mandiri. & Mahasiswa mengerjakan aktivitas terarah untuk backup dan pemulihan sistem dan menunjukkan evidence hasil belajar. & Bentuk penilaian: Backup dan pemulihan sistem. Mengacu rubrik RTM. & Ketepatan konsep; kualitas artefak/praktik; validasi dan dokumentasi. & 5 \\
16 & SCPMK0208-01304 & Tugas besar remastering Linux. & Modalitas: Blended Learning. Bentuk: Praktikum. Metode: demonstrasi, praktik laboratorium Linux, problem based learning, dan diskusi. & 1 x 2 x 50' tatap muka; 1 x 2 x 50' tugas/praktik mandiri. & Mahasiswa mengerjakan aktivitas terarah untuk tugas besar remastering linux dan menunjukkan evidence hasil belajar. & Bentuk penilaian: Tugas besar remastering Linux. Mengacu rubrik RTM. & Ketepatan konsep; kualitas artefak/praktik; validasi dan dokumentasi. & 5 \\
17 & SCPMK0502-01303, SCPMK0208-01304 & UAS validasi administrasi sistem. & Modalitas: Blended Learning. Bentuk: Praktikum. Metode: demonstrasi, praktik laboratorium Linux, problem based learning, dan diskusi. & 1 x 2 x 50' tatap muka; 1 x 2 x 50' tugas/praktik mandiri. & Mahasiswa mengerjakan aktivitas terarah untuk uas validasi administrasi sistem dan menunjukkan evidence hasil belajar. & Bentuk penilaian: UAS validasi administrasi sistem. Mengacu rubrik RTM. & Ketepatan konsep; kualitas artefak/praktik; validasi dan dokumentasi. & 25 \\
\end{longtblr}
